\documentclass[12pt,a4paper]{article}
\usepackage[utf8]{vietnam}
\usepackage{amsmath,amssymb}
\usepackage{tabularx}
\usepackage{tikz}
\usetikzlibrary{calc,angles,quotes,arrows.meta,patterns}
\usepackage{xcolor}
\usepackage{enumitem}
\usepackage{fancyhdr}
\usepackage{graphicx}
\usepackage[top=1.2cm, bottom=1.2cm, left=1.5cm, right=1.5cm]{geometry}
\usepackage[dethi]{ex_test}
\SAOPTN{widthSA=2.6,heightSA=0.62}

\setlength{\headheight}{15pt}
\setlength{\parskip}{0pt}

\pagestyle{fancy}
\fancyhf{}
\rhead{\textit{Môn Toán 6 -- Đề số 15}}
\lhead{\textit{Bộ đề ôn tập Giữa học kì I}}
\cfoot{\thepage}

\begin{document}

% --- KHUNG TIÊU ĐỀ ĐỀ THI ---
\noindent
\begin{tabularx}{\linewidth}{X r}
\textbf{TRƯỜNG THCS: \dotfill} & \textbf{ĐỀ KIỂM TRA ĐỊNH KỲ GIỮA HỌC KỲ I} \\
\textbf{Họ và tên: \dotfill} & \textbf{NĂM HỌC 2026 -- 2027} \\
\textbf{Lớp: \dotfill \quad SBD: \dotfill} & \textbf{Môn: TOÁN 6} \\
\hline
\end{tabularx}
\vspace{2mm}

\begin{center}
  \textbf{\large ĐỀ ÔN TẬP GIỮA HỌC KỲ I -- ĐỀ SỐ 15}\\
  \textit{Thời gian làm bài: 90 phút (Không kể thời gian phát đề)}
\end{center}
\vspace{1mm}

% --- PHẦN I ---
\noindent\textbf{\large PHẦN I. CÂU TRẮC NGHIỆM NHIỀU PHƯƠNG ÁN LỰA CHỌN (3,0 điểm)}

\textit{Thí sinh trả lời từ câu 1 đến câu 12. Mỗi câu hỏi thí sinh chỉ chọn một phương án duy nhất.}

\vspace{1mm}
\begin{ex}
Cho tập hợp $A = \{a; 5; b; 7\}$. Cách viết nào sau đây là đúng?
\choice
{\True $5 \in A$}
{$0 \in A$}
{$7 \notin A$}
{$a \notin A$}
\end{ex}

\begin{ex}
Cho hình vẽ sau:
\begin{center}
\includegraphics[width=0.45\linewidth]{crops/test_p57_xref3497.png}
\end{center}
Trong hình vẽ trên có:
\choice
{2 hình thang, 3 hình lục giác đều, 3 hình tam giác đều}
{2 hình thoi, 3 hình lục giác đều, 2 hình thang cân}
{\True 3 hình lục giác đều, 2 hình thang cân, 2 hình tam giác đều}
{2 hình tam giác đều, 1 hình lục giác đều, 4 hình thoi}
\end{ex}

\begin{ex}
Các số La Mã $\text{XII; XXIII; XXIX}$ được đọc lần lượt là:
\choice
{\True 12; 23; 29}
{12; 23; 28}
{11; 23; 28}
{11; 23; 29}
\end{ex}

\begin{ex}
Cho tập hợp $M = \{1; 2; 3; 4; 5; 6\}$. Khẳng định nào sau đây là đúng?
\choice
{\True $5 \in M$}
{$4 \notin M$}
{$8 \in M$}
{$6 \notin M$}
\end{ex}

\begin{ex}
Trong các khẳng định sau, khẳng định nào sai?
\choice
{$2^3 \cdot 2^4 = 2^7$}
{\True $201^0 = 0$}
{$5^7 : 5^5 = 25$}
{$2^4 = 16$}
\end{ex}

\begin{ex}
Giá trị của biểu thức $22 - 2 \cdot (3^2 + 1)$ bằng:
\choice
{8}
{140}
{200}
{\True 2}
\end{ex}

\begin{ex}
Số nào sau đây là số nguyên tố?
\choice
{6}
{1}
{\True 7}
{0}
\end{ex}

\begin{ex}
$\text{ƯCLN}(8, 20, 4)$ là:
\choice
{1}
{20}
{\True 4}
{8}
\end{ex}

\begin{ex}
Cho tập hợp $A$ gồm các số tự nhiên nhỏ hơn 5. Tập $A$ được viết là:
\choice
{$A = \{1; 2; 3; 4\}$}
{\True $A = \{0; 1; 2; 3; 4\}$}
{$A = \{1; 2; 3; 4; 5\}$}
{$A = \{0; 1; 2; 3; 4; 5\}$}
\end{ex}

\begin{ex}
Cho hình vuông $ABCD$. Khẳng định nào sau đây là sai?
\choice
{Hình vuông $ABCD$ có bốn cạnh bằng nhau: $AB = BC = CD = DA$}
{Hình vuông $ABCD$ có bốn góc ở mỗi đỉnh $A, B, C, D$ bằng nhau}
{Hình vuông $ABCD$ có hai đường chéo bằng nhau: $AC = BD$}
{\True Hình vuông $ABCD$ có hai cặp cạnh đối song song: $AB$ và $BC$; $CD$ và $DA$}
\end{ex}

\begin{ex}
Hình thang cân không có tính chất nào dưới đây?
\choice
{Hai cạnh bên bằng nhau}
{Hai đường chéo bằng nhau}
{Hai cạnh đáy song song}
{\True Hai góc kề một cạnh bên bằng nhau}
\end{ex}

\begin{ex}
Trong các số sau: 2; 7; 11; 21, số nào là hợp số?
\choice
{2}
{7}
{\True 21}
{11}
\end{ex}


\vspace{1mm}
\noindent\textbf{BẢNG TRẢ LỜI TRẮC NGHIỆM PHẦN I:}
\begin{center}
\begin{tabular}{|c|c|c|c|c|c|c|c|c|c|c|c|c|}
\hline
\textbf{Câu} & \textbf{1} & \textbf{2} & \textbf{3} & \textbf{4} & \textbf{5} & \textbf{6} & \textbf{7} & \textbf{8} & \textbf{9} & \textbf{10} & \textbf{11} & \textbf{12} \\
\hline
\textbf{Chọn} & \rule{0pt}{0.65cm} & & & & & & & & & & & \\
\hline
\end{tabular}
\end{center}
\vspace{2mm}

\setcounter{ex}{0}
\noindent\textbf{\large PHẦN II. CÂU TRẮC NGHIỆM ĐÚNG SAI (2,0 điểm)}

\textit{Thí sinh trả lời câu 1 và câu 2. Trong mỗi ý a), b), c), d) ở mỗi câu, thí sinh chọn đúng hoặc sai.}

\vspace{1mm}
\begin{ex}
Cho tập hợp $X = \{x \in \mathbb{N} \mid 5 \le x < 10\}$. Xét tính đúng sai của các mệnh đề sau:
\choiceTF
{\True Tập hợp $X$ gồm 5 phần tử}
{Số 10 là một phần tử thuộc tập hợp $X$}
{Giá trị của biểu thức $3^2 + 4^2$ là một phần tử thuộc tập hợp $X$}
{\True Tổng tất cả các phần tử trong tập hợp $X$ bằng 35}
\end{ex}

\begin{ex}
Cho hình tam giác đều $ABC$ có chu vi bằng $18\text{ cm}$. Xét tính đúng sai của các mệnh đề sau:
\choiceTF
{\True Độ dài mỗi cạnh của tam giác $ABC$ là $6\text{ cm}$}
{\True Ba góc tại các đỉnh $A, B, C$ của tam giác $ABC$ bằng nhau}
{\True Ghép 6 tam giác đều bằng nhau có cạnh $6\text{ cm}$ ta sẽ tạo thành một hình lục giác đều cạnh $6\text{ cm}$}
{Chu vi của hình lục giác đều được tạo thành ở trên là $48\text{ cm}$}
\end{ex}


\noindent\textbf{BẢNG TRẢ LỜI PHẦN II:}
\begin{center}
\begin{tabular}{|c|c|c|c|c||c|c|c|c|c|}
\hline
\textbf{Câu 1} & \textbf{a)} & \textbf{b)} & \textbf{c)} & \textbf{d)} & \textbf{Câu 2} & \textbf{a)} & \textbf{b)} & \textbf{c)} & \textbf{d)} \\
\hline
\textbf{Đ/S} & \rule{0pt}{0.6cm} & & & & \textbf{Đ/S} & & & & \\
\hline
\end{tabular}
\end{center}
\vspace{2mm}

\setcounter{ex}{0}
\noindent\textbf{\large PHẦN III. CÂU HỎI TRẢ LỜI NGẮN (2,0 điểm)}

\textit{Thí sinh trả lời câu 1 và câu 2. Điền kết quả số vào ô trống tương ứng.}

\vspace{1mm}
\begin{ex}
Một trường Trung học cơ sở có 220 học sinh đi tham quan. Nhà trường cần thuê ít nhất bao nhiêu xe ô tô 45 chỗ ngồi để đủ chỗ cho tất cả học sinh? \shortans{5}
\end{ex}

\begin{ex}
Hình chữ nhật có chiều dài và chiều rộng lần lượt là $6\text{ cm}$ và $4\text{ cm}$. Chu vi hình chữ nhật đó là bao nhiêu centimét? \shortans{20}
\end{ex}


\setcounter{ex}{0}
\noindent\textbf{\large PHẦN IV. TỰ LUẬN (3,0 điểm)}

\textit{Thí sinh trình bày chi tiết lời giải vào giấy làm bài.}

\vspace{1mm}
\begin{ex} (2,0 điểm)
Thực hiện các phép tính sau (tính hợp lý nếu có thể):
\begin{enumerate}[label=\textbf{\alph*)}, leftmargin=0.8cm, itemsep=2pt]
\item (0,5 điểm) $270 : 3 : 5 - 4$
\item (0,5 điểm) $9^3 \cdot 2^3 - 4^2 : 2^4$
\item (0,5 điểm) $867 - (167 + 80)$
\item (0,5 điểm) $600 : \{450 : [450 - (4 \cdot 5^3 - 2^3 \cdot 5^2)]\}$
\end{enumerate}
\end{ex}
\vspace{1.5mm}

\begin{ex} (1,5 điểm)
Tìm số tự nhiên $x$, biết:
\begin{enumerate}[label=\textbf{\alph*)}, leftmargin=0.8cm, itemsep=2pt]
\item (0,25 điểm) $x - 15 = 35$
\item (0,25 điểm) $2(x + 15) = 5^2 \cdot 2$
\item (0,5 điểm) $128 - 3(x + 4) = 23$
\item (0,5 điểm) $x \ \vdots \ 15$ và $0 < x \le 60$
\end{enumerate}
\end{ex}
\vspace{1.5mm}

\begin{ex} (1,0 điểm)
Một nền nhà hình chữ nhật của gia đình bác An có chiều dài là $12\text{ m}$, chiều rộng là $8\text{ m}$.
\begin{enumerate}[label=\textbf{\alph*)}, leftmargin=0.8cm, itemsep=2pt]
\item (0,5 điểm) Tính diện tích nền nhà hình chữ nhật đó.
\item (0,5 điểm) Bác An dự định lát nền bằng những viên gạch men hình vuông có cạnh là $40\text{ cm}$. Hỏi bác An phải dùng hết bao nhiêu viên gạch men để lát kín nền nhà đó? (biết diện tích mạch vữa là không đáng kể).
\end{enumerate}
\end{ex}
\vspace{1.5mm}

\begin{ex} (0,5 điểm)
Cho biểu thức $B = 2 + 21^{21} + 25^{120}$. Chứng minh rằng $B$ là hợp số.
\begin{enumerate}[label=\textbf{\alph*)}, leftmargin=0.8cm, itemsep=2pt]
\item (0,5 điểm) Chứng minh $B$ là hợp số.
\end{enumerate}
\end{ex}
\vspace{1.5mm}


\begin{center}
\rule{4cm}{0.4pt}\\
\textit{--- HẾT ---}
\end{center}
\end{document}
